% !TEX root = ../../main.tex

\chapter{卷积神经网络}

在前面的章节中，我们已经学习了神经网络的基本结构、反向传播、参数优化以及训练与泛化等内容，对神经网络如何学习已经有了较为完整的认识。接下来，我们将目光转向图像数据，看看神经网络应该如何处理图像。图像由大量像素组成，这些像素按照二维空间排列，彼此之间还存在明显的空间关系。那么，我们应该如何利用这些信息，使神经网络能够更有效地理解和识别图像中的内容？

\section{全连接的局限性}

我们首先从一个经典的入门任务开始。如\cref{fig:mnist_7}所示，我们希望设计一个神经网络，使其能够识别MNIST数据集\footnote{MNIST（Modified National Institute of Standards and Technology database）是机器学习与计算机视觉领域经典的手写数字图像数据集，由$0$到$9$共10类手写数字组成。该数据集共包含70,000张图像，其中训练集60,000张，测试集10,000张。每张图像的分辨率为$28\times28$，为单通道灰度图像，并配有对应的数字类别标签。由于数据规模适中、任务直观，MNIST常被用于图像分类模型的入门实验。}中的手写数字$0$到$9$。也就是说，给定一张手写数字图像，神经网络需要判断其中的数字属于$0$到$9$中的哪一类。

\begin{figure}[H]
    \boxedimage[width=0.25\textwidth]{figures/part_02/p2_ch01/mnist_7.pdf}
    \caption{MNIST手写数字图片示例（数字7）。}
    \label{fig:mnist_7}
\end{figure}

一张MNIST图片的分辨率为$28\times28$，是一张单通道灰度图像，因此包含$28\times28=784$个像素。为了让神经网络能够处理这张图片，我们首先需要执行\textbf{展平（Flatten）}操作，将原本的二维图像矩阵转换成一个包含$784$个元素的一维向量。随后，我们接入两个隐藏层，神经元数量分别为$512$和$256$，最后接入一个包含$10$个神经元的输出层，分别对应数字$0$到$9$，其网络结构如\cref{fig:fcn}所示。从图中可以看到，该网络中相邻两层之间的连接非常密集：上一层的每一个神经元，都与下一层的所有神经元相连接。因此，这种网络结构也被称为\textbf{全连接网络（Fully Connected Network）}。

接下来，我们计算一下这个网络的参数量。对于一个全连接层，每个输入神经元与每个输出神经元之间都有一个权重，因此权重矩阵包含$N_{\text{in}}\times N_{\text{out}}$个参数；同时，每个输出神经元还对应一个偏置，因此偏置向量包含$N_{\text{out}}$个参数。因此，对于一个输入节点数为$N_{\text{in}}$、输出节点数为$N_{\text{out}}$的全连接层，其参数量为：
$$
\text{参数量} = (N_{\text{in}} \times N_{\text{out}}) + N_{\text{out}}
$$

根据这个公式，我们可以进一步计算该网络各层的参数量：
\begin{itemize}
    \item 输入层 $\to$ 隐藏层1：输入784，输出512。参数量为$784 \times 512 + 512 = 401,920$。
    \item 隐藏层1 $\to$ 隐藏层2：输入512，输出256。参数量为$512 \times 256 + 256 = 131,328$。
    \item 隐藏层2 $\to$ 输出层：输入256，输出10。参数量为$256 \times 10 + 10 = 2,570$。
\end{itemize}

\begin{figure}[H]
    \boxedimage[width=0.95\textwidth]{figures/part_02/p2_ch01/fcn.pdf}
    \caption{全连接网络处理手写数字的结构示意图。左侧$28\times 28$图像被展平为784维输入向量，随后依次送入两个隐藏层（分别含512和256个神经元），最终由10个输出节点给出各类别得分。}
    \label{fig:fcn}
\end{figure}

将各层的参数量相加，可以得到该模型的总参数量：$401,920+131,328+2,570=535,818$。也就是说，仅仅是一个分辨率为$28\times28$、只有$10$个类别的简单图像，就需要超过$53$万个参数。对于MNIST这样的简单任务来说，这样的参数规模似乎还可以接受；但当图像分辨率进一步提高时，参数量会迅速增长。假设我们需要处理一张分辨率为$224\times224$、包含RGB三个通道的彩色图像，那么经过展平之后，输入维度将达到$224\times224\times3=150,528$。如果第一层隐藏层仅设置$1,024$个神经元，那么仅这一层的参数量就达到：
$$
150,528\times1,024+1,024=154,141,696
$$
也就是说，单单一个全连接层就包含约$1.54$亿个参数。如果继续堆叠多个全连接层，参数量还会进一步增加。这就是全连接网络处理图像时面临的第一个问题：\textbf{参数爆炸（Parameter Explosion）}。

庞大的参数量不仅意味着更高的计算和存储成本，也会增加模型训练的负担。但更重要的是，全连接网络在处理图像时，还难以充分利用图像本身的空间结构。如\cref{fig:flatten}所示，为了让全连接网络能够处理图像，我们必须先将二维的像素矩阵展平为一维向量。展平操作虽然没有改变像素本身的数值，但改变了像素原有的组织方式，破坏了图像原有的二维空间结构，像素之间的空间位置信息也不再被显式保留。

在原始图像中，相邻像素之间通常具有较强的局部相关性。以手写数字“7”为例，横杠、斜杠以及它们之间的连接区域，都是由空间上相邻的像素共同构成的，这些局部像素之间的组合关系对于识别数字形状非常重要。然而，图像展平之后，原本的二维位置关系被转换成了一维排列，网络接收到的只是一个长向量，很难直接利用像素之间的局部空间关系。对于全连接层而言，展平后的数字“7”与一串随机排列的像素值在输入形式上并没有区别，网络需要通过训练才能逐渐学习其中的空间规律。这是全连接网络处理图像时面临的第二个问题：\textbf{空间结构信息丢失（Loss of Spatial Structure）}。

\begin{figure}[H]
    \boxedimage[width=0.95\textwidth]{figures/part_02/p2_ch01/flatten.pdf}
    \caption{图像展平操作示意图（为便于可视化，此处以$10\times10$图像为例）。上方为二维像素矩阵，经展平操作后按行首尾相接排列成下方的一维向量；红色竖线标出原图像各行在一维向量中的对应位置。展平后，像素数值保持不变，但原有的二维空间结构不再被显式表示。}
    \label{fig:flatten}
\end{figure}

除了难以利用空间结构之外，全连接网络对图像中的位置变化也较为敏感。假设数字“7”向左上方移动了几个像素。对于人类而言，我们依然能够轻松判断它是数字“7”；但对于全连接网络而言，输入向量中的大量元素都会随之发生变化。由于不同位置的像素对应不同的权重，同一个特征出现在不同位置时，网络接收到的像素组合也会随之发生变化。这意味着，当数字“7”出现在图像中的不同位置时，全连接网络需要针对不同位置重新学习对应的像素组合关系。同一个特征换一个位置，原来的参数关系就难以直接复用，只能重新学习。图像中可能出现的位置越多，需要学习的参数关系也就越复杂。如\cref{fig:translate}所示，数字“7”仅向左上方平移了几个像素，展平后的一维向量中就已经有大量元素发生了变化。这就是全连接网络处理图像时面临的第三个问题：\textbf{平移不变性缺失（Lack of Translation Invariance）}。

\begin{figure}[H]
    \boxedimage[width=0.95\textwidth]{figures/part_02/p2_ch01/translate.pdf}
    \caption{平移导致一维向量发生剧变（为便于可视化，此处以$10\times 10$图像为例）。与\cref{fig:flatten}对照：数字7仅向左上方平移了几个像素，但展平后的一维向量中已有大量元素的取值发生显著改变。}
    \label{fig:translate}
\end{figure}

那么，有没有一种网络结构，既能够直接利用图像的空间结构，又能够用更少的参数有效地提取图像中的局部特征呢？

\section{卷积操作}

让我们先暂时抛开复杂的数学公式，回想一下人类是如何观察图像的。以数字“7”为例，我们通常会关注顶部横杠、斜杠走向和拐角等局部区域，而不会逐一分析图像中的$784$个像素。如\cref{fig:mnist_7_local}所示，这些局部区域包含了识别数字形状的重要信息。既然图像中的特征具有明显的局部性，那么神经网络也不必连接所有像素，而是可以只关注图像中的一个局部区域，从中提取相应的特征。

\begin{figure}[H]
    \boxedimage[width=0.25\textwidth]{figures/part_02/p2_ch01/mnist_7_local.pdf}
    \caption{数字7的局部特征。红色方框标出识别时关注的局部区域，如顶部横杠、拐角与斜杠走向。}
    \label{fig:mnist_7_local}
\end{figure}

图像中可能存在大量这样的局部区域，如果每个区域都使用一组独立的权重来提取特征，参数量仍然很大。但是，无论竖直边缘出现在图像的左上角还是右下角，本质上都是同一种特征，因此我们没有必要为不同位置分别学习参数。我们可以只学习一组参数，并将它应用到图像的不同位置，在整张图像上滑动寻找相同的特征。这样，同一组参数就可以被多个位置重复使用，这就是\textbf{参数共享（Weight Sharing）}。

\begin{figure}[H]
    \boxedimage[width=0.85\textwidth]{figures/part_02/p2_ch01/convolution.pdf}
    \caption{一次卷积操作。红框标出$3\times 3$感受野，其像素值与卷积核（中）对应位置的权重逐元素相乘并求和，得到右侧单个特征响应值。}
    \label{fig:once_convolution_}
\end{figure}

如\cref{fig:once_convolution_}所示，在输入图像上滑动的局部窗口称为\textbf{感受野（Receptive Field）}。与之对应相乘的权重矩阵则被称为\textbf{卷积核（Kernel）}，也叫\textbf{滤波器（Filter）}。在实际计算中，卷积核按照预设的\textbf{步长（Stride）}在输入图像上滑动，步长表示卷积核每次移动的像素数。例如，步长为$1$时，卷积核每次移动$1$个像素；步长为$2$时，每次移动$2$个像素。卷积核每移动到一个位置，就将感受野内的像素值与卷积核中对应位置的权重逐元素相乘并求和，得到该位置的特征响应值，这一计算过程就是\textbf{卷积操作（Convolution）}。

\begin{figure}[H]
    \boxedimage[width=0.75\textwidth]{figures/part_02/p2_ch01/full_convolution.pdf}
    \caption{完整的卷积操作（步长为1）。卷积核自左上角起逐位置滑动，红框为当前感受野；每滑到一个位置便计算一个响应值填入右侧特征图，直至遍历全图、得到完整的$8\times 8$输出。}
    \label{fig:full_convolution}
\end{figure}

卷积核遍历整张图像后，每个位置都会得到一个特征响应值，这些响应值按照原来的空间位置排列起来，便形成一张新的二维矩阵，即\textbf{特征图（Feature Map）}。完整的滑动过程如\cref{fig:full_convolution}所示。以图中的卷积核为例，它主要用于检测水平线，因此遍历整张图像后，水平线区域会产生较强的响应，最终在特征图中突出显示数字“7”顶部的横杠，而其他区域的响应相对较弱。

但一张图像中并不只有水平线，还可能包含竖直线、斜线、圆弧等各种不同的局部特征。因此，在实际的卷积神经网络中，我们通常不会只使用一个卷积核，而是使用多个不同的卷积核，每个卷积核负责提取一种或一类局部特征。例如，一个卷积核可能更容易响应水平边缘，另一个卷积核可能更容易响应竖直边缘，还有一些卷积核则可以响应更加复杂的局部纹理。每个卷积核都会在整张图像上完整地滑动一次，并生成一张对应的特征图。多个卷积核得到的特征图堆叠在一起，便构成卷积层的输出。多卷积核的卷积过程如\cref{fig:multi_channel_convolution}所示。

\begin{figure}[H]
    \boxedimage[width=0.95\textwidth]{figures/part_02/p2_ch01/multi_channel_convolution.pdf}
    \caption{多通道卷积操作。输入为6个通道的特征图（$10\times 10\times 6$），卷积核相应地含6个通道片（$3\times 3\times 6$），红框为跨通道的感受野。计算时先在每个通道内做$3\times 3$卷积、各得一个标量（下方逐通道展开），再将6个标量跨通道求和（加偏置），得到输出特征图（$8\times 8$）上的单个响应值。}
    \label{fig:multi_channel_convolution}
\end{figure}

这里需要特别注意，图中的$3\times3\times6$表示的是一个卷积核，其中前两个维度$3\times3$表示卷积核在空间上的大小，最后一个维度$6$表示它需要同时覆盖输入的$6$个通道。因此，这个卷积核会对6个通道分别进行卷积，再将各通道的结果相加，最终得到一张$8\times8$的单通道特征图。也就是说，卷积核的通道数必须与输入特征图的通道数一致，而一个卷积核只能产生一个输出通道。如果希望卷积层输出多个通道，就需要使用多个这样的卷积核。如\cref{fig:multi_channel_output}所示，使用$3$个$3\times3\times6$的卷积核，就可以得到$3$张$8\times8$的特征图，最终形成一个$8\times8\times3$的多通道输出。

\begin{figure}[H]
    \boxedimage[width=0.75\textwidth]{figures/part_02/p2_ch01/multi_channel_output.pdf}
    \caption{多卷积核的多通道卷积操作。输入为$10\times 10\times 6$的特征图；3个彼此独立的$3\times 3\times 6$卷积核各自完成一次跨通道卷积，分别产出一张$8\times 8$的单通道特征图；3张特征图堆叠在一起，最终形成$8\times 8\times 3$的多通道输出。}
    \label{fig:multi_channel_output}
\end{figure}

理解了卷积核如何提取特征之后，我们再用数学公式描述这一过程。假设输入图像为二维矩阵$X$，卷积核为大小为$K\times K$的二维矩阵$W$，则卷积核在位置$(i,j)$处的输出为：
$$
Y(i,j)=\sum_{m=0}^{K-1}\sum_{n=0}^{K-1}X(i+m,j+n)\,W(m,n).
$$
这个公式表示：将卷积核$W$放在输入图像的位置$(i,j)$上，把感受野中的像素值与卷积核中的对应位置的权重逐元素相乘，再将所有乘积相加，得到输出$Y$在该位置的特征响应值。随后，卷积核按照设定的步长移动到下一个位置，并重复这一过程，直到遍历整个图像。

不过，我们可以发现一个问题：当卷积核只在输入图像内部滑动，而不允许超出边界时，输出的空间尺寸会逐渐缩小。如\cref{fig:full_convolution}所示，对于$10\times10$的输入图像，使用$3\times3$的卷积核、步长$s=1$且不进行填充时，输出特征图的大小变为$8\times8$。如果继续堆叠多层卷积，这种尺寸缩小会不断累积，使特征图越来越小，从而限制网络的深度。与此同时，边缘像素能够参与卷积计算的次数也少于位于中央的像素，这意味着图像边缘区域中的信息更容易受到影响。

\begin{figure}[H]
    \boxedimage[width=0.95\textwidth]{figures/part_02/p2_ch01/padding.pdf}
    \caption{零填充示意图。为防止特征图尺寸缩小，在输入图像边缘填充零值，从而保持特征图尺寸不变。}
    \label{fig:padding}
\end{figure}

为了解决这一问题，我们可以在输入图像的四周补充一些额外的像素。最常见的方式是填充数值为$0$的像素，因此称为\textbf{零填充（Zero Padding）}。填充的厚度，也就是在输入图像中每一侧增加的像素层数，记为$p$。如\cref{fig:padding}所示，图中在原图像的上下左右各增加了$1$圈$0$像素，因此$p=1$。通过在图像边缘增加填充，卷积核可以覆盖原本位于边缘的区域，同时还能控制输出特征图的空间尺寸。

一般情况下，对于输入高度为$H_{\text{in}}$、卷积核大小为$K$、填充为$p$，步长为$s$的卷积操作，输出高度为：
$$
H_{\text{out}} = \left\lfloor \frac{H_{\text{in}}+2p-K}{s} \right\rfloor+1
$$
其中$\lfloor \cdot \rfloor$表示向下取整，即取不超过该数的最大整数。宽度的计算方式与之相同。以图中的设置为例，输入图像大小为$10\times10$，卷积核大小为$3\times3$，步长$s=1$，填充$p=1$，因此输出高度为：
$$
H_{\text{out}} = \left\lfloor \frac{10+2\times1-3}{1} \right\rfloor+1 =10
$$
输出宽度同样为$10$，因此最终得到的特征图尺寸仍为$10\times10$，实现了输入与输出空间尺寸保持不变。

到这里，我们描述的主要是卷积层中的加权求和过程。实际上，一个卷积神经元的完整计算还包括偏置和激活函数。首先，在加权求和的结果上加上一个\textbf{偏置}。通常，一个卷积核对应一个独立的偏置参数$b$，并将其加到该卷积核在不同位置产生的响应值上。其次，还需要通过\textbf{激活函数}引入非线性。如果没有激活函数，多层卷积的堆叠在数学上仍然可以表示为线性变换，网络就难以拟合复杂的非线性模式。因此，在实际的卷积神经网络中，卷积层后通常会接ReLU等激活函数。

对于单通道输入，一个卷积核在位置$(i,j)$处的完整计算过程可以写成：
$$
Y(i,j) = \operatorname{ReLU} \left( \sum_{m=0}^{K-1} \sum_{n=0}^{K-1} X(i+m,j+n)\cdot W(m,n)+b \right)
$$
其中，$W$表示卷积核的权重，$b$表示偏置，$\operatorname{ReLU}$表示激活函数。当输入包含多个通道时，例如输入特征图的尺寸为$H\times W\times C_{\text{in}}$，一个卷积核也需要同时覆盖全部$C_{\text{in}}$个通道，因此其尺寸为$K\times K\times C_{\text{in}}$。此时，卷积核会分别对每个输入通道中的局部区域进行加权求和，再将所有通道的结果相加，最后加上偏置并通过激活函数。因此，多通道情况下，在位置$(i,j)$处的计算可以表示为：
$$
Y(i,j) = \operatorname{ReLU} \left( \sum_{c=0}^{C_{\text{in}}-1} \sum_{m=0}^{K-1} \sum_{n=0}^{K-1} X(i+m,j+n,c)\cdot W(m,n,c) +b \right)
$$
其中，$c$表示输入通道的索引。一个$K\times K\times C_{\text{in}}$的卷积核会综合所有输入通道的信息，最终得到一个输出通道中的一个响应值。如果使用多个卷积核，那么每个卷积核都会生成一张独立的特征图，所有特征图堆叠起来，就构成卷积层的多通道输出。

可以看到，无论是单通道还是多通道，卷积操作的核心仍然是对局部区域进行加权求和，再加上偏置并通过激活函数。这与全连接神经元的计算形式
$$
y=f\left(\sum_i w_i x_i+b\right)
$$
在结构上是一致的。不同之处在于，全连接层中的一个神经元会连接所有输入，而卷积层只连接局部区域；同时，同一个卷积核的权重会在不同位置重复使用，从而实现局部连接和权重共享。

\section{池化操作}

有了卷积之后，我们可以通过堆叠多个卷积层不断提取图像中的局部特征。但如果每一层都通过零填充保持特征图的空间尺寸不变，那么随着网络加深，较大的特征图会使后续层始终需要处理大量数据，从而带来更高的计算开销。同时，较高的空间分辨率也保留了更加精细的位置信息。例如，数字“7”的顶部横线稍微移动几个像素，卷积核虽然仍然能够检测到它，但特征图中的响应位置也会随之变化，使模型对这种局部位置变化更加敏感。那么，有没有一种操作，既能够降低特征图的空间尺寸，又能够适当弱化特征对精确位置的依赖呢？

\begin{figure}[H]
    \boxedimage[width=0.75\textwidth]{figures/part_02/p2_ch01/pooling.pdf}
    \caption{最大池化操作示意图（$2\times 2$窗口、步长2）。红框为当前池化窗口，窗口内像素的最大值被填入右侧输出特征图的对应位置；窗口以步长2滑遍全图后，$10\times 10$的输入被下采样为$5\times 5$的输出。}
    \label{fig:pooling}
\end{figure}

我们在识别数字“7”时，并不需要知道一条横线精确位于第几行第几列。只要能够确认图像上半部分存在一条横线，我们就可以利用这一信息继续判断数字的形状。因此，我们可以在一个局部区域内保留最具有代表性的特征，同时舍弃部分精确的位置信息。这个操作就是\textbf{池化（Pooling）}，它是一种常见的\textbf{下采样（Downsampling）}操作。一个直接的做法是，在局部区域内只保留响应最强的特征，其余信息直接舍弃，这就是\textbf{最大池化（Max Pooling）}。如\cref{fig:pooling}所示，我们使用一个$2\times2$的窗口，以步长$2$在特征图上滑动，每次只保留窗口中的最大值。这样，特征图的空间尺寸可以直接缩小一半，而通道数保持不变。

为什么只保留最大值？因为经过卷积和激活后，较大的响应通常意味着对应位置的特征更加明显。例如，当卷积核用于检测水平边缘时，包含明显水平边缘的区域往往会产生较大的响应，因此最大池化可以保留局部区域中最显著的特征。与此同时，池化缩小了特征图的空间尺寸，减少了后续卷积层需要处理的数据量，从而降低计算开销；而且只要特征仍然位于同一个池化窗口内，即使发生轻微移动，最大值也可能保持不变，因此能够减弱模型对特征精确位置的依赖，使其对小范围的位置变化更加稳定。

除了最大池化之外，还有\textbf{平均池化（Average Pooling）}，它计算池化窗口内所有元素的平均值。与最大池化相比，平均池化会综合考虑窗口中的所有响应，因此能够保留区域的整体信息，但如果一个局部区域中只有少数位置具有较强的特征响应，而其他位置的响应较弱，那么较大的响应就可能被周围较小的值平均，从而使显著特征被削弱。如\cref{fig:pooling_comparison}所示，最大池化能够保留局部区域中最明显的特征，而平均池化则使这些特征被周围较弱的响应“稀释”。

\begin{figure}[H]
    \boxedimage[width=0.75\textwidth]{figures/part_02/p2_ch01/pooling_comparison.pdf}
    \caption{最大池化与平均池化的对比（均为$2\times 2$窗口、步长$2$）。最大池化取每个窗口内的最大值，以保留局部最强响应；平均池化则取窗口内所有元素的平均值，因而输出更加平滑。}
    \label{fig:pooling_comparison}
\end{figure}

到这里，我们已经学习了卷积和池化这两个基本操作。在实际的卷积神经网络中，它们通常会按照“卷积$\rightarrow$激活$\rightarrow$池化”的方式反复堆叠，形成多层网络。随着网络逐层加深，不同层学习到的特征也会逐渐变得更加复杂：浅层网络直接处理原始像素，主要关注边缘、横线、竖线、拐角等基础特征；中间层在这些基础特征的基础上进一步组合，逐渐识别出弧线、局部形状和更加复杂的几何结构；深层网络则继续整合这些局部特征，形成更加抽象的整体结构，为最终的图像识别提供更丰富的特征表示。也就是说，卷积神经网络通过层层堆叠，让特征从简单的像素和边缘逐步发展为更加抽象的高级特征。

\section{全连接}

假设我们处理一张$28\times28$大小的手写数字图像，经过若干层卷积和池化后，最终得到一个大小为$5\times5$、通道数为$16$的特征图。此时，这$16$个通道中的数值已经不再是原始的像素灰度值，而是不同卷积核提取出的特征在各个位置上的响应。接下来，我们需要利用这些特征判断输入图像究竟属于数字$0$到$9$中的哪一类。前面的卷积和池化主要负责从图像中提取和压缩特征，使网络逐步得到更加抽象的特征表示，但这些操作本身并不直接完成最终的类别判断。我们还需要一种能够综合不同特征之间关系的结构，将前面提取出的特征进一步组合，并映射到具体的类别上，这就是全连接层所要完成的任务。

我们首先将$5\times5\times16$的特征图展平（Flatten）成一个$400$维的一维向量，然后接入一层或多层全连接层，对这些特征进行组合，最后通过一个包含$10$个神经元的输出层，分别对应数字$0$到$9$。输出层得到的$10$个数称为\textbf{logits}，它们表示模型对各个类别的相对判断结果，但本身并不是概率，可以是任意实数。接下来，我们对这$10$个logits进行Softmax运算：
$$
p_i=\frac{e^{z_i}}{\sum_{j=1}^{K}e^{z_j}}
$$
这个运算称为\textbf{Softmax函数}。其中，$z_i$表示第$i$个类别的logit，$K$表示类别总数，在MNIST任务中$K=10$。经过Softmax之后得到的$p_i$均为$0$到$1$之间的数，并且所有类别对应的$p_i$之和恰好等于$1$，因此可以将其解释为模型预测输入图像属于第$i$类的概率。

Softmax主要完成两个步骤。首先，通过指数函数$e^{z_i}$将任意实数转换为正数，同时拉开不同logit之间的差距，使较大的logit对应更大的数值。随后，将每个指数值除以所有类别指数值之和进行归一化，从而使所有输出之和等于$1$。因此，logit越大的类别，经过Softmax后通常会获得越高的预测概率。例如，假设网络对一张图片输出的$10$个logits中，最大的三个分别为$z_7=5.0$、$z_1=2.0$和$z_9=1.0$，其余$7$个类别的logit都近似为$0$。那么数字$7$的预测概率约为：
$$
p_7=\frac{e^{5.0}}{e^{5.0}+e^{2.0}+e^{1.0}+7e^0}\approx\frac{148.4}{148.4+7.4+2.7+7}\approx0.90.
$$
因此，模型会以约$90\%$的概率判断这张图片是数字$7$，而数字$1$和数字$9$的概率分别约为$4.5\%$和$1.6\%$，其余类别的概率则更低。

至此，我们就得到了一个完整的卷积神经网络：\redbf{卷积负责提取局部特征，池化负责压缩空间尺寸并减弱对精确位置的依赖，多层堆叠逐步形成更加抽象的特征，全连接层负责综合这些特征并完成分类，最后通过Softmax得到各类别的预测概率}。

\section{LeNet-5}

我们已经从零开始构建出了一个完整的卷积神经网络。实际上，这套思路早在20世纪80年代末至90年代初就已经被提出并应用。1989年，Yann LeCun等人发表了\href{http://yann.lecun.com/exdb/publis/pdf/lecun-89e.pdf}{\textit{Backpropagation Applied to Handwritten Zip Code Recognition}}，将局部连接、权重共享和下采样等思想应用于手写邮政编码识别，这是早期卷积网络的重要工作。此后，LeCun等人在\href{http://yann.lecun.com/exdb/publis/pdf/lecun-90c.pdf}{\textit{Handwritten Digit Recognition with a Back-Propagation Network}}等工作中继续发展这一架构，并将其应用于手写字符识别。1998年，Yann LeCun、Léon Bottou、Yoshua Bengio和Patrick Haffner在\href{http://yann.lecun.com/exdb/publis/pdf/lecun-98.pdf}{\textit{Gradient-Based Learning Applied to Document Recognition}}中系统阐述了以\textbf{LeNet-5}为代表的卷积神经网络\footnote{LeNet中的“Le”取自Yann LeCun的姓氏，“Net”是network的缩写，“5”则表示它是LeNet系列的第五个版本。相比此前的版本，LeNet-5的结构更加成熟，在MNIST测试集上识别准确率约为99.05\%，因此成为LeNet系列中最具代表性的版本。}，并展示了其在文档识别任务中的效果。LeNet-5是早期卷积神经网络中最具代表性的经典架构，奠定了现代CNN的基本结构框架，对后续卷积神经网络的发展产生了深远影响。其经典架构如\cref{fig:lenet}所示。

\begin{figure}[H]
    \boxedimage[width=0.95\textwidth]{figures/part_02/p2_ch01/lenet.pdf}
    \caption{LeNet-5网络结构示意图。$32\times 32$输入经“卷积—下采样”交替堆叠（C1/S2、C3/S4）逐步提取特征，末端由全连接层C5（120）、F6（84）映射到10个类别输出。}
    \label{fig:lenet}
\end{figure}

LeNet-5的具体流程如下：
\begin{enumerate}
    \item \textbf{输入层}：接收一张$32\times32$的单通道灰度图像，即$32\times32\times1$。需要注意的是，MNIST原始图像大小为$28\times28$，LeNet-5将其扩展到$32\times32$，在数字周围保留一定的边缘空间，从而避免数字过于靠近图像边界，为后续卷积提供更充分的空间。
    \item \textbf{卷积层C1}：使用$6$个$5\times5$卷积核，步长$s=1$，不进行填充$p=0$。因此，空间尺寸由$32$变为$32-5+1=28$，通道数由$1$变为$6$，输出为$28\times28\times6$。
    \item \textbf{池化层S2}：采用$2\times2$平均池化，步长$s=2$，使空间尺寸减半，输出为$14\times14\times6$。
    \item \textbf{卷积层C3}：使用$16$个$5\times5$卷积核，步长$s=1$，不进行填充$p=0$。空间尺寸由$14$变为$14-5+1=10$，输出为$10\times10\times16$。
    \item \textbf{池化层S4}：采用$2\times2$平均池化，步长$s=2$，空间尺寸再次减半，输出为$5\times5\times16$。
    \item \textbf{卷积层C5}：此时输入特征图的空间尺寸为$5\times5$，而C5使用的卷积核同样为$5\times5$。因此，每个卷积核在空间上只能覆盖整个特征图并计算一次，不会产生多个空间位置的响应。从计算结果来看，这相当于将$5\times5\times16=400$个输入同时连接到一个神经元，因此C5在功能上等价于全连接层。C5共使用$120$个卷积核，因此得到$120$个输出神经元。
    \item \textbf{全连接层F6}：将$120$个神经元映射到$84$个神经元。
    \item \textbf{输出层}：将$84$个神经元映射到$10$个输出，分别对应数字$0$到$9$。
\end{enumerate}

现在，让我们计算一下LeNet-5的总参数量：

\begin{enumerate}
    \item \textbf{卷积层C1}：共有$6$个$5\times5\times1$的卷积核，每个卷积核包含$5\times5\times1+1=26$个参数，其中$1$个为偏置，因此参数量为$6\times26=156$。
    \item \textbf{卷积层C3}：共有$16$个$5\times5\times6$的卷积核，每个卷积核包含$5\times5\times6+1=151$个参数，因此参数量为$16\times151=2,416$。
    \item \textbf{池化层S2、S4}：没有可学习的参数，因此参数量为$0$。
    \item \textbf{卷积层C5}：使用$120$个$5\times5\times16$的卷积核，每个卷积核包含$5\times5\times16+1=401$个参数（含偏置），因此参数量为$120\times401=48,120$。
    \item \textbf{全连接层F6}：将$120$个神经元连接到$84$个神经元，因此参数量为$120\times84+84=10,164$。
    \item \textbf{输出层}：将$84$个神经元连接到$10$个输出神经元，因此参数量为$84\times10+10=850$。
\end{enumerate}

将以上各层的参数量相加，可以得到：
$$
156+2,416+48,120+10,164+850=61,706.
$$
也就是说，按照这种连接方式计算，整个网络只需要约$6.2$万个参数。作为对比，前面介绍的全连接网络$784\rightarrow512\rightarrow256\rightarrow10$共有$535,818$个参数。可以看到，通过局部连接和参数共享，卷积神经网络在大幅减少参数量的同时，仍然能够逐层提取图像中的特征。

\begin{quote}
    至此，我们已经完整搭建了卷积神经网络的架构，并认识了其中的核心组件。那么，卷积神经网络又该如何进行训练呢？
\end{quote}